\documentclass[11pt]{article}
\usepackage[margin=1in]{geometry}
\usepackage{amsmath,amssymb,amsthm,mathtools}
\usepackage{graphicx}
\usepackage[hidelinks]{hyperref}
\usepackage{microtype}

\newtheorem{theorem}{Theorem}
\newtheorem{lemma}{Lemma}
\newtheorem{remark}{Remark}
\newcommand{\R}{\mathbb R}
\newcommand{\T}{\mathbb T}
\newcommand{\eps}{\varepsilon}
\newcommand{\vol}{\operatorname{vol}}

\title{Two minimizing log-polynomials in every even degree at least eight}
\author{Candidate counterexample for arXiv:2007.07952\\
\small autonomous proof-search run \texttt{fa\_banach\_001}}
\date{13 August 2026}

\begin{document}
\maketitle

\begin{center}
\fbox{\parbox{0.91\linewidth}{\textbf{Status.} Candidate full negative
answer, likely valid, pending expert review.  For every even $d\geq8$ we
construct a normalized integrable log-concave function on $\R^2$ for which
Problem~1 of the source paper has exactly two minimizers.}}
\end{center}

\section{Source question}

Alonso-Guti\'errez, Gonz\'alez Merino and Villa consider, for even $d$, the
problem of minimizing
\begin{equation}\label{eq:problem}
 \int_{\R^n}t\exp\{-g(x)^{1/d}\}\,dx
   =t n!\,|G_1(g)|,
 \qquad G_1(g)=\{x:g(x)\leq1\},
\end{equation}
over $t\geq1$ and nonnegative homogeneous degree-$d$ polynomials $g$, subject
to
\[
 f(x)\leq t\exp\{-g(x)^{1/d}\}.
\]
They prove existence for normalized integrable log-concave $f$.  On PDF
page~4, after Theorem~1.1, they note that uniqueness is known for $d=2$ but is
not known in general, and that they know no example with a nonunique
minimization point \cite{AGMV}.

\begin{figure}[ht]
\centering
\includegraphics[width=.94\linewidth]{figures/open_problem_crop.png}
\caption{The residual uniqueness question in the source paper (PDF page~4).}
\end{figure}

We give such examples for every even $d\geq8$.

\section{Proof intuition}

Write $f=e^{-u}$ and put $r=\log t$.  The source reduction says that the
fixed-$r$ polynomial must contain
\[
 H_{e^r}(f)=\{u^*\leq r\}^{\circ}.
\]
Thus we may design the sublevel bodies $D_r=\{u^*\leq r\}$ first and then
recover $f$ by convex duality.

We take a nested family $D_r$ whose middle piece linearly interpolates, by
Minkowski addition, between two scaled and rotated polynomial dual balls.
At the endpoints, the optimal outer polynomial bodies are known exactly.
Along the middle piece, symmetry forces every optimizer to have only the two
Fourier modes $1$ and $\cos(d\theta)$.  Its containment constraint becomes the
problem of finding the largest affine minorant of a concave power mean, hence
the optimizer is exactly the chord joining the endpoint values.

The endpoint scale is selected so that their weighted volumes agree.  A
uniform Taylor calculation then gives
\[
 \log\frac{e^{r_\lambda}v(r_\lambda)}{e^{r_0}v(r_0)}
 =\frac2d\left(1-\frac6d\right)\lambda(1-\lambda)\eps^2
  +O\bigl(\lambda(1-\lambda)\eps^3\bigr).
\]
For $d\geq8$ this is positive throughout $0<\lambda<1$ when $\eps$ is small.
Simple homotheties make the objective strictly decreasing before the first
endpoint and strictly increasing after the second.  The endpoint pairs are
therefore the two global minimizers.

\section{Reduction to a family of polar bodies}

For a compact set $K\subset\R^2$ containing the origin, let
\[
 L_d(K)=\min\left\{\,|G_1(g)|\ \middle|\
 \begin{array}{l}
 g\text{ is a nonnegative homogeneous degree-$d$ polynomial,}\\[-2pt]
 K\subseteq G_1(g)
 \end{array}\right\}.
\]
Lasserre's theorem gives a unique optimizing polynomial \cite{Lasserre}; we
use the theorem in precisely the form recorded in the source paper.

\begin{lemma}[Realizing nested bodies]\label{lem:realize}
Let $(D_r)_{r\geq0}$ be compact origin-symmetric convex bodies satisfying
\begin{enumerate}
\item $D_a\subseteq D_b$ for $a\leq b$;
\item $D_{(1-\lambda)a+\lambda b}\supseteq
 (1-\lambda)D_a+\lambda D_b$;
\item $D_0=\{0\}$ and $\bigcup_{r\geq0}D_r=\R^2$.
\end{enumerate}
Then there is an even finite convex function $w$ with $w(0)=0$ and
$\{w\leq r\}=D_r$.  If $u=w^*$ and $f=e^{-u}$, then $f$ is even,
log-concave, $f(0)=\|f\|_\infty=1$, integrable, and
\begin{equation}\label{eq:polar-reduction}
 \overline{H_{e^r}(f)}=D_r^\circ .
\end{equation}
Consequently the best fixed-$r$ value in \eqref{eq:problem}, with the fixed
factor $2!$ omitted, is
\begin{equation}\label{eq:objective}
 \Phi(r)=e^r L_d(D_r^\circ).
\end{equation}
\end{lemma}

\begin{proof}
Set $w(x)=\inf\{r:x\in D_r\}$.  The nesting gives
$\{w\leq r\}=D_r$ (closedness also handles the infimum), and the Minkowski
concavity gives convexity of $w$.  The remaining assertions except
integrability follow from elementary Fenchel duality.  In particular
$u^*=w$, and the polarity calculation in Lemma~3.3 of \cite{AGMV} gives
\eqref{eq:polar-reduction}.

In our application below, $D_r=(r/r_0)D_+$ for $0\leq r\leq r_0$.  Thus
$w(x)=r_0\|x\|_{D_+}$ wherever $w(x)\leq r_0$, and
\[
 u(y)=\sup_x(\langle x,y\rangle-w(x))
 \geq \sup_{0\leq a\leq1}a(h_{D_+}(y)-r_0)
 =(h_{D_+}(y)-r_0)_+.
\]
Since $D_+$ contains a Euclidean ball, $h_{D_+}(y)\geq c|y|$ for some
$c>0$.  Hence $e^{-u(y)}\leq e^{-(c|y|-r_0)_+}$ is integrable.
The equivalence between feasibility and $H_{e^r}(f)\subseteq G_1(g)$ proved
in Lemma~3.1 of \cite{AGMV} now yields \eqref{eq:objective}.
\end{proof}

\section{The construction and the exact fixed-height optimizer}

Fix an even integer $d\geq8$, write $\alpha=1/d$, and fix $\kappa>0$.  For a
sufficiently small $\eps\in(0,1)$, use polar coordinates
$x=\rho(\cos\theta,\sin\theta)$
and define
\[
 p_+(x)^d=\rho^d(1+\eps\cos d\theta),\qquad
 p_-(x)^d=s^d\rho^d(1-\eps\cos d\theta),
 \qquad \sigma=\left(\frac{1+\eps}{1-\eps}\right)^\alpha,\qquad
 s=\sigma e^{\kappa\eps^2}.
\]
The expressions on the right are homogeneous polynomials: they equal
$(x^2+y^2)^{d/2}\pm\eps\operatorname{Re}(x+iy)^d$, with the second one
multiplied by $s^d$.  For small $\eps$, $p_\pm$ are norms.  Indeed the unit
circle of $p_+$ has radial function
$R(\theta)=(1+\eps\cos d\theta)^{-1/d}$, and
$R^2+2(R')^2-RR''=1+O_d(\eps)$ uniformly; the same applies to $p_-$.  Let
$C_\pm=\{p_\pm\leq1\}$ and $D_\pm=C_\pm^\circ$.

The ratio $p_-/p_+$ ranges from $s/\sigma=e^{\kappa\eps^2}$ to $s\sigma$.
Thus $C_-\subset\operatorname{int}C_+$ and, after polarity,
$D_+\subset\operatorname{int}D_-$.  Put
\[
 \Delta=2\log s,\qquad r_0=\frac12,\qquad r_1=r_0+\Delta,
 \qquad
 \eta=\frac{(s/\sigma)-1}{2\Delta(s/\sigma)}.
\]
and define, for $r\geq0$,
\begin{equation}\label{eq:Dr}
D_r=
\begin{cases}
 \dfrac{r}{r_0}D_+,&0\leq r\leq r_0,\\
 (1-\lambda)D_++\lambda D_-,&r=r_0+\lambda\Delta,
       \quad0\leq\lambda\leq1,\\
 \bigl(1+\eta(r-r_1)\bigr)D_-,&r\geq r_1.
\end{cases}
\end{equation}
This family is increasing.  Its support function is concave in $r$: it is
linear on each of the three intervals, and at the two knots its slopes obey
\[
 \frac{h_{D_-}-h_{D_+}}{\Delta}\leq\frac{h_{D_+}}{r_0},
 \qquad
 \frac{h_{D_-}-h_{D_+}}{\Delta}\geq\eta h_{D_-}.
\]
Indeed, on writing $q=h_{D_-}/h_{D_+}=p_-/p_+$, we have
$s/\sigma\leq q\leq s\sigma$.  The first assertion follows for small
$\eps$ from
\[
 \frac{s\sigma-1}{\Delta}\longrightarrow1<\frac1{r_0}=2,
\]
and the second follows directly from
$(q-1)/(\Delta q)\geq((s/\sigma)-1)/(\Delta(s/\sigma))=2\eta$.
Moreover $\eta>0$, so the family exhausts $\R^2$, and $\eta\to0$ as
$\eps\to0$.  Hence Lemma~\ref{lem:realize} applies.

We next compute $L_d(D_r^\circ)$ on the middle interval.  Since the support
function of $D_\pm=C_\pm^\circ$ is $p_\pm$, the polar of the Minkowski sum in
\eqref{eq:Dr} has radial function
\begin{equation}\label{eq:radial}
 R_\lambda(\theta)
 =\frac{1}{(1-\lambda)(1+\eps\cos d\theta)^\alpha
       +\lambda s(1-\eps\cos d\theta)^\alpha}.
\end{equation}

\begin{lemma}[Exact optimizer on the middle segment]\label{lem:optimizer}
For $0\leq\lambda\leq1$, set
\begin{align*}
 a_\lambda&=(1-\lambda)(1+\eps)^\alpha
                 +\lambda s(1-\eps)^\alpha,\\
 b_\lambda&=(1-\lambda)(1-\eps)^\alpha
                 +\lambda s(1+\eps)^\alpha,\\
 A_\lambda&=\frac{a_\lambda^d+b_\lambda^d}{2},\qquad
 B_\lambda=\frac{a_\lambda^d-b_\lambda^d}{2}.
\end{align*}
Then the unique Lasserre polynomial of $D_{r_0+\lambda\Delta}^\circ$ is
\begin{equation}\label{eq:g-lambda}
 g_\lambda(\rho,\theta)=\rho^d(A_\lambda+B_\lambda\cos d\theta).
\end{equation}
In particular
\begin{equation}\label{eq:volume}
 L_d(D_{r_0+\lambda\Delta}^\circ)
 =\frac12\int_0^{2\pi}
 (A_\lambda+B_\lambda\cos d\theta)^{-2/d}\,d\theta.
\end{equation}
\end{lemma}

\begin{proof}
The body in \eqref{eq:radial} is invariant under reflection and rotation by
$2\pi/d$.  Uniqueness of the Lasserre polynomial makes the optimizer invariant
under these symmetries.  A real homogeneous degree-$d$ polynomial restricted
to the circle has Fourier frequencies $d,d-2,\ldots,-d$.  Rotation by
$2\pi/d$ retains only $0,\pm d$, and reflection removes the sine term; hence
it has the form $\rho^d(A+B\cos d\theta)$.

Containment is equivalent, with $z=\cos d\theta$, to
\begin{equation}\label{eq:minorant}
 A+Bz\leq
 F_\lambda(z):=\bigl((1-\lambda)(1+\eps z)^\alpha
       +\lambda s(1-\eps z)^\alpha\bigr)^d,\qquad -1\leq z\leq1.
\end{equation}
The weighted power mean $M_\alpha(x,y)=((1-\lambda)x^\alpha+\lambda
y^\alpha)^{1/\alpha}$ is concave and coordinatewise increasing for
$0<\alpha\leq1$ (the reverse Minkowski inequality).  Since
$x=1+\eps z$ and $y=s^d(1-\eps z)$ are affine in $z$, $F_\lambda$ is
concave.  Therefore its endpoint chord
\[
 \ell_\lambda(z)=\frac{F_\lambda(1)+F_\lambda(-1)}2
 +\frac{F_\lambda(1)-F_\lambda(-1)}2z
 =A_\lambda+B_\lambda z
\]
is feasible.  Every other feasible affine function $q(z)$ obeys
$q(\pm1)\leq\ell_\lambda(\pm1)$, so the affine difference
$q-\ell_\lambda$ is nonpositive throughout $[-1,1]$.  Thus its polynomial
sublevel body contains that of \eqref{eq:g-lambda}, proving optimality.
Formula \eqref{eq:volume} is the polar-coordinate area formula.
\end{proof}

\section{The two minima}

\begin{lemma}[Middle barrier]\label{lem:barrier}
For every fixed even $d\geq8$, all sufficiently small $\eps>0$ satisfy
\[
 e^{r_0+\lambda\Delta}L_d(D_{r_0+\lambda\Delta}^\circ)
 >e^{r_0}L_d(C_+)\qquad(0<\lambda<1),
\]
whereas equality holds for $\lambda=0,1$.
\end{lemma}

\begin{proof}
Let $V_\lambda$ denote the right-hand side of \eqref{eq:volume} and define
\[
 Q(\lambda,\eps)=\lambda\Delta+\log V_\lambda-\log V_0.
\]
At $\lambda=0$, $g_\lambda=p_+^d$; at $\lambda=1$,
$g_\lambda=p_-^d$.  Rotation does not change area, so
$V_1=s^{-2}V_0$ and $Q(0,\eps)=Q(1,\eps)=0$ exactly.

All quantities are real analytic near the compact square
$[0,1]\times\{0\}$.  Direct Taylor expansion, using
\[
 \frac1{2\pi}\int_0^{2\pi}\cos(d\theta)\,d\theta=0,
 \qquad
 \frac1{2\pi}\int_0^{2\pi}\cos^2(d\theta)\,d\theta=\frac12,
\]
gives, uniformly for $0\leq\lambda\leq1$,
\begin{equation}\label{eq:taylor}
 Q(\lambda,\eps)
 =2\alpha(1-6\alpha)\lambda(1-\lambda)\eps^2
   +O\bigl(\lambda(1-\lambda)\eps^3\bigr).
\end{equation}
For completeness, here is an auditable coefficient calculation.  Uniformly
in $\lambda$,
\begin{align*}
 A_\lambda&=1+2\lambda\eps+a_2(\lambda)\eps^2+O(\eps^3),\\
 B_\lambda&=(1-2\lambda)\eps+b_2(\lambda)\eps^2+O(\eps^3),\\
 a_2(\lambda)&=\lambda\left(4\alpha(1-\lambda)+4\lambda-2
                    +\frac{\kappa}{\alpha}\right),\\
 b_2(\lambda)&=2\lambda(2\alpha\lambda-2\alpha-2\lambda+1),\\
 \Delta&=4\alpha\eps+2\kappa\eps^2+O(\eps^3).
\end{align*}
Writing $c_\lambda=B_\lambda/A_\lambda$, the angular average in
\eqref{eq:volume} gives
\[
 \log V_\lambda
 =\log\pi-2\alpha\log A_\lambda
  +\frac{\alpha(2\alpha+1)}2c_\lambda^2+O(\eps^3).
\]
Substitution and subtraction of the value at $\lambda=0$ cancel both the
linear term and the $\kappa$ term, leaving
\[
 2\alpha(1-6\alpha)\lambda(1-\lambda).
\]
The displayed remainder has the endpoint factors because $Q$ is analytic and
vanishes identically at $\lambda=0,1$: successive analytic division by
$\lambda$ and $1-\lambda$ makes the quotient bounded on a fixed
neighborhood.  Since $1-6\alpha>0$ for $d\geq8$, \eqref{eq:taylor} is
positive for every $0<\lambda<1$ once $\eps$ is sufficiently small.
\end{proof}

\begin{theorem}[Negative answer to the uniqueness question]\label{thm:main}
For every even integer $d\geq8$ there is an even normalized integrable
log-concave function $f:\R^2\to[0,1]$ for which Problem~1 of \cite{AGMV}
has exactly two minimizing pairs.
\end{theorem}

\begin{proof}
Choose $\eps$ as in Lemma~\ref{lem:barrier}, small enough also that
$\eta<1/2$, construct $D_r$ by \eqref{eq:Dr}, and let $f$ be supplied by
Lemma~\ref{lem:realize}.  On $0<r\leq r_0$, homogeneity gives
\[
 L_d(D_r^\circ)=L_d((r_0/r)C_+)=(r_0/r)^2L_d(C_+),
\]
so $\Phi(r)$ is strictly decreasing there because $r\leq r_0=1/2<2$.
On $r\geq r_1$,
\[
 L_d(D_r^\circ)=\bigl(1+\eta(r-r_1)\bigr)^{-2}L_d(C_-),
\]
and
\[
 \frac{d}{dr}\log\Phi(r)
 =1-\frac{2\eta}{1+\eta(r-r_1)}>0.
\]

The endpoint separation is $\Delta=2\log s$, so
$e^{r_1}|C_-|=e^{r_0+\Delta}s^{-2}|C_+|=e^{r_0}|C_+|$.
Here $L_d(C_\pm)=|C_\pm|$, because $C_\pm$ are themselves the sublevel
bodies of $p_\pm^d$, while no containing body can have smaller volume.
Lemma~\ref{lem:barrier} puts every point of the middle segment strictly above
this common value.  The two outer pieces
move away from the common value, so the only global minimizers occur at
$r=r_0$ with polynomial $p_+^d$ and at $r=r_1$ with polynomial $p_-^d$.
They are distinct because $\eps>0$.
\end{proof}

\section{Verification and limitations}

The accompanying script independently checks the symbolic coefficient in
\eqref{eq:taylor} and samples the exact integral \eqref{eq:volume}.  Those
checks are sanity tests, not proof.  The argument proves nonuniqueness for
every even $d\geq8$ in dimension two.  It does not decide the borderline
degrees $d=4$ and $d=6$: the quadratic coefficient is negative for $d=4$ and
zero for $d=6$.

A bounded novelty search on 13 August 2026 covered the run indexes, exact
title/id, author/title combinations, and arXiv/web searches for
log-polynomial, $d$-Lasserre--L\"owner, uniqueness, and nonunique minimizers.
It found no later resolution.  This is not proof of novelty.

\paragraph{Human-review recommendation.}
Check (i) realization of the nested family, (ii) the support-slope bookkeeping
at the two knots, (iii) the invariant-polynomial reduction, (iv) concavity and
chord optimality in Lemma~\ref{lem:optimizer}, and (v) the uniform endpoint
factor in \eqref{eq:taylor}.

\begin{thebibliography}{9}
\bibitem{AGMV}
D. Alonso-Guti\'errez, B. Gonz\'alez Merino, and R. Villa,
\emph{Best approximation of functions by log-polynomials},
J. Funct. Anal. 280 (2021), 108999; arXiv:2007.07952.

\bibitem{Lasserre}
J.-B. Lasserre,
\emph{A generalization of L\"owner--John's ellipsoid theorem},
Math. Program. 152 (2015), 559--591; arXiv:1302.1056.
\end{thebibliography}

\end{document}
